\section{The Experiment Index}

Every claim in this paper rests on a committed, deterministic experiment. Nothing is asserted on authority. Where the body says that a self heals its own damage, or that the field ground state obeys an area law, or that the matter knot is a fermion, there is a named program in the project's test tree that was run, that varied size rather than seeds, and that reported a verdict either way. This appendix is the one place those experiment ids belong, gathered into clean tables so that a reader can trace any statement in the paper back to the run that grounds it.

The registry holds $457$ experiments across $16$ categories. Each is a single experiment for a single claim. Each is fully deterministic, which is a load-bearing constraint and not a convenience. The base of the theory is discrete and reproducible, so a result that depended on a random seed would not be a result about the theory at all. To measure at a larger scale, an experiment grows its mesh, not its randomness. This is the deepest reason the index can serve as evidence rather than anecdote.

\subsection{The depth rubric}

Each experiment carries a depth level. The level is an honesty grade for what the experiment actually proves, and it is stated up front so that no result is read for more than it is worth. The rubric is owned in full by the section on the depth rubric and detailed in the methodology appendix. In brief there are four levels.

\begin{definition}[The four depth levels]
An experiment is graded \textbf{L0} if it is circular or a baseline control, \textbf{L1} if it establishes known mathematics on the substrate, \textbf{L2} if it reproduces known physics on the substrate, and \textbf{L3} if its result is emergent and novel and carries a control that could have failed.
\end{definition}

The grade L3 is the demanding one. An L3 experiment is not allowed to put its conclusion in by hand. It must run a control in which the effect could have come out absent, and it must report the contrast rather than only the success. The eight base atoms (mesh, dock, site, tone, lean, knit, wake, beat) are fixed before any L3 experiment runs, so an emergent result cannot have been smuggled into the base for the occasion.

\subsection{The registry by depth}

The depth profile of the whole registry is shown in Table~\ref{tab:depth-registry}. It is deliberately honest about its own shape. The large middle of the registry reproduces known mathematics and known physics on the substrate, which is the necessary groundwork. The load-bearing claims are the $44$ at level L3.

\begin{table}[ht]
\centering
\begin{tabular}{@{}lr@{}}
\toprule
depth & count \\
\midrule
L3, emergent with a falsifiable control & 44 \\
L2, known physics reproduced & 284 \\
L1, known mathematics established & 114 \\
L0, baseline or control & 15 \\
\midrule
\textbf{total} & \textbf{457} \\
\bottomrule
\end{tabular}
\caption{The registry graded by depth.}
\label{tab:depth-registry}
\end{table}

\begin{principle}[The shape is the honesty]
Most experiments reproduce known mathematics or known physics on the substrate. The $44$ at level L3 are the load-bearing results, and each one carries a control that could have failed and reports the outcome either way. A registry weighted the other way, all L3 and no groundwork, would be the suspicious one.
\end{principle}

\subsection{The registry by category}

The same $457$ experiments sort into $16$ categories by what they probe, from the selves that the theory is built to explain down to the addressing machinery that lets a dock be named. Table~\ref{tab:category-registry} gives the count in each category and, in the last column, how many of those reach level L3. The two heaviest concentrations of L3 results sit in selves and in holography, which is where the theory's strongest novel claims live.

\begin{table}[ht]
\centering
\begin{tabular}{@{}lrr@{}}
\toprule
category & count & L3 of those \\
\midrule
selves & 114 & 18 \\
gauge & 39 & 1 \\
spin & 31 & 0 \\
relativity & 31 & 2 \\
substrate-survey & 30 & 1 \\
data-structure & 28 & 0 \\
cosmology & 27 & 5 \\
foundations & 25 & 0 \\
geometry & 23 & 1 \\
quantum & 22 & 1 \\
gravity & 21 & 0 \\
holography & 17 & 8 \\
renormalization & 16 & 2 \\
associative & 16 & 5 \\
computation & 10 & 0 \\
addressing & 7 & 0 \\
\bottomrule
\end{tabular}
\caption{The registry graded by category, with the L3 count of each.}
\label{tab:category-registry}
\end{table}

\subsection{The strongest results, every L3 experiment}

The $44$ emergent results are the ones a skeptic should read first, because each could have failed and did not. They are listed here in full, grouped by category, with the verdict each returned. An experiment id is a path in the test tree, so each row points at a runnable program.

\subsubsection{Selves}

The selves category carries the most L3 results, $18$ of them, and they are the heart of the theory's claim that a persistent, self-correcting pattern can emerge from the eight atoms and maintain itself with no outside knower. Table~\ref{tab:l3-selves} lists them.

\begin{table}[ht]
\centering
\begin{tabular}{@{}ll@{}}
\toprule
id & verdict \\
\midrule
selves/autonomous-self & a self maintains itself by a purely local knit, no outside knower \\
selves/emergent-self-robust & a self emerges by cohesion, conserving maintenance holds it against decay \\
selves/horosphere-self & selves are far more persistent on the flat horosphere than in the bulk \\
selves/integrated-information & selves are tone-integration local maxima \\
selves/memory-vs-conservation & charge is conserved while the pattern decays \\
selves/metacognition & the self-model predicts the self's next state \\
selves/p121-recursion & a hub represents a lower hub, recursion of representation \\
selves/p152-evolution & heredity plus variation plus selection drives mean fitness up \\
selves/p153-integrated-agent & the closed perceive-plan-act loop crosses a sequence of barriers \\
selves/p161-evolution & selection raises fitness \\
selves/p162-integrated-agent & multi-step lookahead solves a detour the reactive agent cannot \\
selves/permanent-memory & a maintained codeword stays at full fidelity where unmaintained erodes \\
selves/persistent-self & a self-maintaining region keeps its identity and boundary \\
selves/self-emergence & fixed fills do not self-organize selves, the honest negative \\
selves/self-maintenance & a self heals its own damage by the knit alone \\
selves/self-model & a localized hub represents the self's global state \\
selves/selves-at-scale & coherent self-patches emerge on the exact substrate \\
selves/selves-interacting & opposite selves annihilate at contact \\
\bottomrule
\end{tabular}
\caption{The $18$ L3 experiments in the selves category.}
\label{tab:l3-selves}
\end{table}

\subsubsection{Holography}

Holography carries $8$ L3 results, the densest L3 concentration relative to its size, and they are the strongest support for the claim that the boundary of the mesh encodes its interior. Table~\ref{tab:l3-holography} lists them.

\begin{table}[ht]
\centering
\begin{tabular}{@{}ll@{}}
\toprule
id & verdict \\
\midrule
holography/area-law & the field ground state is area-law while a thermal state is volume-law \\
holography/black-hole & region entropy scales with horizon area, not volume \\
holography/bulk-nonlocality & distant surface points are joined by a short hidden path through the bulk \\
holography/growing-code & the growing holographic code raises its erasure threshold with shell age \\
holography/holographic-memory & a spread-encoded bit survives a bounded erasure, a localized blob is destroyed \\
holography/p91-holography & the Ryu-Takayanagi log law \\
holography/ryu-takayanagi-73 & the area law on the drawable two-dimensional face \\
holography/signaling & a signal crosses the universe through the bulk to a far self \\
\bottomrule
\end{tabular}
\caption{The $8$ L3 experiments in the holography category.}
\label{tab:l3-holography}
\end{table}

\subsubsection{Cosmology and associative}

Cosmology and the associative category each carry $5$ L3 results. The cosmology results concern how the growing mesh gives an arrow, an asymmetry, and a rare thin tail of matter. The associative results concern how hyperbolic geometry makes memory search logarithmic where a flat lattice is polynomial. Table~\ref{tab:l3-cosmo-assoc} lists both groups.

\begin{table}[ht]
\centering
\begin{tabular}{@{}ll@{}}
\toprule
id & verdict \\
\midrule
cosmology/baryogenesis & an emergent matter-antimatter asymmetry \\
cosmology/dimension-selection & only three dimensions give stable closed orbits \\
cosmology/growth-expansion & net-positive birth gives emergent expansion, with a static control \\
cosmology/large-n-crossing & a cluster move traverses heights \\
cosmology/rarity-measures & three independent measures confirm life is rare \\
associative/capacity-vs-curvature & capacity rises and search latency falls with curvature \\
associative/hopfield-emergent-recall & recall works on the dissipative layer, fails on the bare reversible knit \\
associative/parallel-cost & search cost is constant in size, the communication radius grows logarithmically \\
associative/search-latency & bulk search latency scales logarithmically with size \\
associative/spreading-activation & spreading activation through semantic memory is the bulk query wave \\
\bottomrule
\end{tabular}
\caption{The $5$ L3 experiments in cosmology and the $5$ in the associative category.}
\label{tab:l3-cosmo-assoc}
\end{table}

\subsubsection{Relativity, renormalization, and the singletons}

The remaining L3 results are spread thin across categories that carry only one or two each. The two relativity results and the two renormalization results sit alongside four singletons, one apiece in gauge, geometry, quantum, and the substrate survey. Table~\ref{tab:l3-rest} lists them. The substrate-survey result is worth singling out, since it is the one that protects the whole program from the most dangerous objection, that a lattice theory is ruled out by gamma-ray-burst timing. Every hyperbolic substrate passes that Lorentz bound where the flat lattice does not.

\begin{table}[ht]
\centering
\begin{tabular}{@{}ll@{}}
\toprule
id & verdict \\
\midrule
relativity/deterministic-wave & a deterministic reversible knit propagates ballistically, so momentum emerges \\
relativity/predictions-vs-bounds & the model passes the gamma-ray-burst Lorentz bound that excludes a lattice \\
renormalization/emergent-macro-rule & a macro-rule emerges on tone-independent blocks, fails when frustrated \\
renormalization/renormalization-keystone & parameters on a small slice match those of a much larger field \\
gauge/mass-hierarchy & unfitted geometric spacing gives a multi-decade mass hierarchy beating a power law \\
geometry/horosphere-flat & a Busemann level set of the honeycomb is the flat cusp \\
quantum/entanglement-bell & the exchange dynamics violates the Bell inequality at the Tsirelson bound \\
substrate-survey/non-random-substrates & every hyperbolic substrate is Lorentz-safe while the flat lattice is not \\
\bottomrule
\end{tabular}
\caption{The remaining L3 experiments, in relativity, renormalization, and four singleton categories.}
\label{tab:l3-rest}
\end{table}

\subsection{The experiment behind each major claim}

The strongest claims of the base and the self trace to a small set of grounding experiments. This section gathers them by the part of the argument they support, so that the chain from claim to run is visible at a glance. The grouping follows the shape of the theory itself, from the base knit and the bath, through the proven negatives, the self-completing additions, the three faces of the self, the decisive discrete findings, and the emergent and nesting results.

\subsubsection{The base knit and the bath}

These experiments ground the lowest layer, that the knit (the collide-then-stream law) persists structure on the substrate, that the open boundary acts as a bath, that the wake supplies the arrow of time, and that a closed system cannot host a self at all. The obstruction result is the sharp one. It says the bath is not optional decoration but a precondition for selfhood. Table~\ref{tab:base-knit} lists them.

\begin{table}[ht]
\centering
\begin{tabular}{@{}p{0.42\textwidth}p{0.32\textwidth}p{0.18\textwidth}@{}}
\toprule
claim & experiment id & verdict \\
\midrule
the bare knit persists structure on the substrate & selves/bare-rule-persistence-3434 & grounds the base knit \\
the open boundary is the bath & selves/bath-from-open-boundary & the bath model \\
the wake gives the arrow and irreversibility & selves/growth-arrow-irreversibility & the arrow from the wake \\
a closed system cannot host a self & selves/substrate-self-obstruction & the obstruction theorem \\
the lone-charge mobility on the sites & selves/mobile-rule-d4, selves/scatter-d4 & mobility and scattering \\
\bottomrule
\end{tabular}
\caption{The experiments grounding the base knit and the bath.}
\label{tab:base-knit}
\end{table}

\subsubsection{The proven negatives}

A theory is trusted more for what it admits it cannot do than for what it claims. These five experiments are honest negatives. They establish that the eight base atoms alone, with nothing added, do not give a bound body. Co-motion is free streaming and not binding. The create move binds but seals the radiation channel. The gas has no restoring pressure. A discrete kink shatters rather than holding. Each of these is a PASS in the sense that the experiment confirmed the negative cleanly. Table~\ref{tab:negatives} lists them.

\begin{table}[ht]
\centering
\begin{tabular}{@{}p{0.42\textwidth}p{0.32\textwidth}p{0.18\textwidth}@{}}
\toprule
claim & experiment id & verdict \\
\midrule
co-motion is free streaming, not binding & selves/co-motion-not-bound & honest negative \\
the create move binds but seals the radiation channel & selves/arrow-binds-but-seals & negative \\
the gas has no restoring pressure & selves/no-restoring-pressure & negative \\
a bound body does not emerge from the eight base atoms alone & selves/no-emergent-bound-body & negative \\
a discrete kink is unstable & selves/discrete-kink-unstable & the kink route fails \\
\bottomrule
\end{tabular}
\caption{The proven negatives, what the eight atoms alone cannot do.}
\label{tab:negatives}
\end{table}

\subsubsection{The self-completing additions}

Because the bound body does not come for free, the theory adds the attraction as an emergent mechanism, binding by radiation, the shadow pressure of the active vacuum. These experiments establish that the addition is minimal and principled rather than arbitrary. The home gives a persistent identity and an energy scale. Vacuum depletion binds at range two with no new field. A minimal bounded field, not arbitrary integers, suffices. A one-to-three-trit ternary field binds and repairs. Table~\ref{tab:additions} lists them, including the negative that bulk curvature disperses rather than binding, which rules out the lazy alternative.

\begin{table}[ht]
\centering
\begin{tabular}{@{}p{0.42\textwidth}p{0.32\textwidth}p{0.18\textwidth}@{}}
\toprule
claim & experiment id & verdict \\
\midrule
the home gives persistent identity and the energy scale & selves/rest-slot-identity & pass \\
active-vacuum depletion binds at range 2, no new field & selves/vacuum-depletion-attraction & pass \\
the bare depletion is sub-critical at long range & selves/casimir-vacuum-attraction & grounds the range limit \\
a minimal bounded field, not arbitrary integers & selves/minimal-attraction-field & pass \\
a 1-to-3-trit ternary field binds and repairs & selves/ternary-field-attraction & pass \\
a full discrete self with identity, repair, and radiation & selves/gravity-bound-self, selves/full-self-positive-control & pass \\
bulk curvature disperses, does not bind & selves/bulk-curvature-disperses & negative \\
\bottomrule
\end{tabular}
\caption{The self-completing additions, with the curvature negative that rules out the alternative.}
\label{tab:additions}
\end{table}

\subsubsection{The three faces of the self}

A self has three faces, identity, agency, and binding, and each face has its own grounding experiment. Identity is a winding reversibly protected at the $24$-site resolution. Agency is a soft mode that emerges with a linear dispersion $\omega = c\,k$, damped by the bath to give the healing channel. Binding is a stabilized soliton that is reversibly bound, with the twist read as a pure rotation carrying no real coupling. The binding face also carries the Derrick result, that pure exchange is scale-marginal so a stabilizer is needed. Table~\ref{tab:three-faces} lists the six.

\begin{table}[ht]
\centering
\begin{tabular}{@{}p{0.12\textwidth}p{0.34\textwidth}p{0.30\textwidth}p{0.10\textwidth}@{}}
\toprule
face & claim & experiment id & verdict \\
\midrule
identity & a winding is reversibly protected at the 24-cell resolution & selves/topological-winding-identity & pass \\
agency & a soft mode emerges, $\omega = c\,k$ & selves/emergent-soft-radiation & pass \\
agency & the bath damps the soft mode, the healing channel & selves/bath-damps-soft-mode & pass \\
binding & a stabilized soliton is reversibly bound & selves/dm-skyrmion-bound-self & pass \\
binding & the twist as pure rotation, no real coupling & selves/quaternion-twist-binding & pass \\
binding & pure exchange is scale-marginal, a stabilizer is needed & selves/derrick-stabilizer & the Derrick result \\
\bottomrule
\end{tabular}
\caption{The three faces of the self and the experiment grounding each.}
\label{tab:three-faces}
\end{table}

\subsubsection{The decisive discrete-self findings}

These four are the sharpest results on the discrete self. A three-trit spin holds the soliton charge exactly. Reversible dynamics conserves the charge at small steps. No finite direction group is fine enough, so the binding dynamics is forced to be emergent rather than baked into a finer lattice. Radiation pressure alone is sub-critical. Together they pin down why the attraction must be emergent and not a base atom. Table~\ref{tab:discrete} lists them.

\begin{table}[ht]
\centering
\begin{tabular}{@{}p{0.42\textwidth}p{0.32\textwidth}p{0.18\textwidth}@{}}
\toprule
claim & experiment id & verdict \\
\midrule
a 3-trit spin holds the soliton charge exactly & selves/ternary-skyrmion-charge & pass \\
reversible dynamics conserves the charge at small steps & selves/reversible-dynamics-step-threshold & pass \\
no finite direction group is fine enough, binding is forced emergent & selves/fine-group-too-coarse & pass \\
radiation pressure alone is sub-critical & selves/reversible-radiation-pressure & the reversible-attraction limit \\
\bottomrule
\end{tabular}
\caption{The decisive discrete-self findings.}
\label{tab:discrete}
\end{table}

\subsubsection{The emergent and nesting results}

The last group lifts the self into its effective description. Selves act as attractors, the L2 effective picture. Selves nest, giving a self of selves. Matter and photon live in one tone field, the one-tone result. Table~\ref{tab:nesting} lists them.

\begin{table}[ht]
\centering
\begin{tabular}{@{}p{0.42\textwidth}p{0.32\textwidth}p{0.18\textwidth}@{}}
\toprule
claim & experiment id & verdict \\
\midrule
selves act as attractors & selves/selves-as-attractors & the L2 effective picture \\
selves nest, a self of selves & selves/nested-selves, selves/nested-bath-selves & nesting \\
matter and photon in one field & selves/leaky-confiner & the one-tone result \\
\bottomrule
\end{tabular}
\caption{The emergent and nesting results.}
\label{tab:nesting}
\end{table}

\subsection{How to read this index}

In the body of the paper an experiment is referred to by a short plain description, the g-factor experiment, the decoupling control, the area-law run. The id is only for cross-reference here. Each id names a deterministic program runnable from the project's test tree, one experiment per claim.

\begin{result}[Grow the size, never the seed]
To measure any of these claims at a larger scale, grow the mesh size and rerun. Never change a seed, because there is no seed to change. The base is deterministic, so the only knob is scale, and a result that holds across scales is a result about the theory rather than about a lucky draw.
\end{result}

The architecture this index catalogs, one experiment per claim with a verdict, is laid out in the section on the experimental architecture. The depth rubric that grades each row is owned by the section on the depth rubric and detailed in the methodology appendix. The base whose claims these experiments ground, the eight atoms and the five pillars they group into, is anatomized in the base anatomy appendix.
